\chapter{Концепція AURA Runtime-орієнтованої інформаційної технології}
\label{ch:concept}

% ============================================================
%  РОЗДІЛ 2 -- концепція, уніфікована модель загроз, межі тверджень.
%  Містить дві опорні таблиці всієї роботи:
%   (1) відображення «проблема -> метод»;
%   (2) межі тверджень (claim boundaries).
%  Ці таблиці роблять докторську розповідь когерентною з початку.
% ============================================================

\section{AURA Runtime як центральний рівень прийняття рішень}
\label{sec:aura-core-idea}
% -- Теза роботи: приватний захист учасника має бути контекстним, stateful і
%    runtime-орієнтованим; криптографія створює умови, за яких він можливий
%    без серверного доступу до відкритого тексту.
% -- AURA Runtime: вхід (подія повідомлення + контекст) -> risk signal -> safety-дія.

\section{Загальна архітектура AURA/Ecliptix}
\label{sec:architecture}
% -- Рівні: (а) автентифікація (PQ-OPAQUE); (б) захищений обмiн (гібридний PQ E2EE);
%    (в) контекстна безпека (AURA Runtime); (г) серверно-сліпий канал нагляду;
%    (д) реалізаційний рівень (Rust services + iOS + FFI + відтворюваність).
% -- Криптографічні рівні -- захисний контур для runtime-рішень.
% [Рисунок: загальна архітектура -- TODO figures/arch_overview]

\section{Потік даних: від повідомлення до safety-рішення}
\label{sec:dataflow}
% -- Локальна/приватна обробка на endpoint; що НЕ залишає пристрій;
%    які межі перетинає лише opaque payload.

\section{Уніфікована модель загроз}
\label{sec:threat-model}
% -- Класичний та квантовий супротивник.
% -- Server-compromise vs endpoint-compromise.
% -- Межа відкритого тексту E2EE; межа витоку метаданих.
% -- Межа операційного керування ключами.
% -- Припущення про довіру до endpoint (integrity) і їхні наслідки.
Запропонована технологія спирається на єдину модель загроз, спільну для всіх рівнів захисту, що дає змогу узгоджено формулювати, які властивості підтримуються криптографічними протоколами, які -- залежать від цілісності endpoint, а які -- підтверджуються експериментальною та release-gated перевіркою. \emph{[Розкрити прозою.]}

\section{Відображення «проблема -- метод»}
\label{sec:problem-method-map}

% Опорна таблиця №1. Робить зв'язок проблем і методів явним.
\begin{longtable}{p{6.0cm} p{6.0cm} p{1.6cm}}
\caption{Відображення проблем приватного контекстного захисту на запропоновані методи}
\label{tab:problem-method}\\
\toprule
\textbf{Проблема} & \textbf{Метод / рівень захисту} & \textbf{Розділ} \\
\midrule
\endfirsthead
\toprule
\textbf{Проблема} & \textbf{Метод / рівень захисту} & \textbf{Розділ} \\
\midrule
\endhead
Передавання пароля або password-equivalent матеріалу серверу & постквантово-посилена OPAQUE-автентифікація & 5 \\
Майбутнє квантове розшифрування перехоплених повідомлень («store-now-decrypt-later») & гібридне постквантове наскрізне шифрування & 6 \\
Дилема safety за наскрізного шифрування (серверна модерація відкритого тексту суперечить приватності) & AURA Runtime: контекстна оцінка ризику + policy-aware рішення & 3--4 \\
Позаконтекстна класифікація окремого повідомлення дає слабкі рішення & метод контекстної (stateful) оцінки інформаційного ризику & 3 \\
Контрольований перегляд (guardian/review) без розкриття відкритого тексту серверу & метод серверно-сліпого каналу нагляду & 7 \\
Доказовість, підзвітність і відтворюваність якості захисту & інформаційна технологія + release-gated оцінювання та privacy-safe аудит & 8 \\
\bottomrule
\end{longtable}

\section{Межі тверджень (claim boundaries)}
\label{sec:claim-boundaries}

% Опорна таблиця №2. Чесно розмежовує доказовість кожного твердження.
% Формулювання ProVerif узгоджено з честю: 3 первинні + 1 scoped honest-message
% correspondence; груповий scoped-модель -- 6 запитів; stress-зобов'язання -- НЕ твердження.
\begin{longtable}{p{6.6cm} p{5.4cm} p{1.6cm}}
\caption{Межі тверджень дисертації за рівнем доказовості}
\label{tab:claim-boundary}\\
\toprule
\textbf{Твердження / властивість} & \textbf{Статус доказовості} & \textbf{Розділ} \\
\midrule
\endfirsthead
\toprule
\textbf{Твердження / властивість} & \textbf{Статус доказовості} & \textbf{Розділ} \\
\midrule
\endhead
Пароль не передається серверу у відкритому або password-equivalent вигляді & доведено формально (Tamarin) & 5 \\
Стійкість до офлайнового перебору словника за компрометації БД без \texttt{oprf\_seed} (класична модель) & доведено формально, з явними операційними межами & 5 \\
Постквантовий захист узгодження сесії в межах гібридної конструкції & доведено формально (scoped) & 5--6 \\
Секретність і кореспонденція активної моделі E2EE & підтверджено формально (ProVerif, scoped): 3 первинні твердження + 1 додаткова scoped-перевірка чесного повідомлення & 6 \\
Серверно-сліпість каналу нагляду (сервер не читає pair key і safety payload) & дизайн + тести/відтворюваний E2E-артефакт & 7 \\
Якість контекстних safety-рішень & оцінено експериментально, release-gated на curated/synthetic корпусах & 3--4, 8 \\
Інженерна реалізованість та інтеграція технології & програмний артефакт + відтворювані перевірки & 8 \\
\textbf{НЕ стверджується:} клінічна/психологічна валідність, абсолютний захист, постквантова стійкість підписів (Ed25519 -- класичний) & поза межами тверджень & --- \\
\bottomrule
\end{longtable}

\rozdilconcl{2}
\emph{[Висновки: сформульована концепція, модель загроз і межі тверджень як спільна рамка для розділів 3--8.]}
